\chapter{Limitações, ética e conclusão}
\label{cap:conclusao}

% Rascunho em tópicos (2026-10-22).
% 2026-10-29 (parecer de orientação): identificadores fora do rosto, origem dos vídeos,
% CEP, fail-open, escopo do custo; E4-lite, política por tamanho e parrot como futuros.

\section{Limitações}
\begin{itemize}
\item Três vídeos de demonstração públicos na internet (29,1\,s); não é um conjunto consentido de crianças.
\item Hiperparâmetros escolhidos nos mesmos vídeos (Seção~\ref{sec:ameacas-validade}).
\item Atacante ingênuo: sem \emph{parrot attack} (reconhecedor retreinado nos rostos anonimizados) \citep{hanisch}.
\item Identificadores fora do rosto não tratados: cabelo, contorno da cabeça, orelhas, roupa, lugar na sala, corpo, voz.
\item Sem avaliação humana de analisabilidade pedagógica.
\item Sem análise de equidade demográfica.
\item Olhar (\emph{gaze}) não medido.
\item Rostos nunca detectados não são ocultados; taxa residual depende de anotação (pendente).
\item Saídas sem troca registradas como geradas (\emph{fail-open}): corrigido em 2026-10-29; frequência nas execuções existentes não medida.
\item Dois reconhecedores; não é garantia formal.
\item Custo: \result{e5/pooled/adapt-0.08/class-hours} GPU-horas por hora de aula, projeção só do laço por rosto; tempo ponta a ponta não medido; só \emph{offline}.
\item Componentes apenas para pesquisa não comercial (Apêndice~\ref{ap:licencas}).
\end{itemize}

\section{Ética e LGPD}
\begin{itemize}
\item Rosto: dado biométrico; etnia: dado sensível; menores: consentimento de um dos pais \citep{lgpd}.
\item Resultado: desidentificação facial, não dado anonimizado (art.~12); o vídeo processado continua dado pessoal.
\item Embeddings reais só em memória; scratch da ponte apagado também em SIGTERM e SIGHUP; crops de rostos não substituídos não são gravados (2026-10-29).
\item Resíduos com fundo real (crops gerados, imagens de identidade, vídeos compostos): removidos com \texttt{experiments/purge.sh} após a exportação.
\item Nenhum quadro de pessoa real no texto; exemplos qualitativos só com rostos sintéticos ou adultos consentidos.
\item Imagens de identidade do SDXL não contam como sintéticas para publicação (cosseno \result{e3g/pooled/identity-cos.plain} com a pessoa real).
\item Vídeos: públicos na internet; fonte, licença e termos por vídeo pendentes de registro.
\item Pendente: consulta ao CEP da UFPR (via orientador e coordenação). Plano B: gravações próprias com adultos voluntários e termo de uso de imagem.
\item Base legal a validar com o orientador e o CEP (não afirmada): LGPD art.~7º, IV e art.~11, II, c, tratamento por órgão de pesquisa para estudos, garantida a anonimização sempre que possível \citep{lgpd}. Esta tese processou biometria de pessoas que não consentiram; o registro da fonte de cada vídeo e a consulta ao CEP são condição para a defesa.
\item Cópias locais dos vídeos e resultados fora do servidor: apagadas (2026-10-29).
\end{itemize}

\section{Trabalhos futuros}
\begin{itemize}
\item Rodar no servidor o que já está implementado: rank inverso, rank-1 por track, rank-1 entre rostos mostrados, auditoria de identidades, $\kappa$ de orientação, estratos por tamanho, Holm; conferir à mão os pares de tracks candidatos e as caixas extras do YOLO-FaceV2-s.
\item Braço $\beta \approx c$ (ortogonal): remove a componente real sem ir para a antípoda.
\item E4-lite, banco de identidades: identidades geradas por prompt genérico, sem condicionar na imagem real e sem inferir atributos; remove a rota de vazamento da imagem de identidade e o custo de difusão por track (1 a 2 dias).
\item E4, identidade só por texto:
  \begin{itemize}
  \item atributos grosseiros (faixa etária, gênero aparente, óculos, cabelo) $\to$ template $\to$ texto-para-imagem, sem ver o rosto real \citep{lee2026,deepprivacy};
  \item remove a rota de vazamento da imagem de identidade (cosseno \result{e3g/pooled/identity-cos.plain}); a troca continua vazando forma;
  \item sem tom de pele ou etnia no vocabulário (LGPD art.~11); validar com compliance e comitê de ética;
  \item riscos: viés e homogeneização do texto-para-imagem, memorização, reidentificação pelo próprio texto.
  \end{itemize}
\item Política por tamanho do rosto: $\beta$ maior, \emph{gate} mais estrito ou ocultação para rostos $\ge 96$\,px (professor).
\item Atacante ciente simplificado (\emph{parrot}: galeria também substituída) \citep{hanisch}.
\item Testes de equivalência (TOST) com margens no piso de ruído.
\item Avaliação em vídeos novos, separados, com o pipeline congelado.
\item Custo da identidade: offload por modelo no lugar do sequencial, sem refiner, LCM-LoRA em 4--5 passos \citep{lcmlora}, DeepCache \citep{deepcache}.
\item Custo por rosto: TensorRT fp16, lotes por quadro, vários workers por GPU, divisão do vídeo entre GPUs; tempo ponta a ponta por estágio.
\item P1 no refiner, ou sem refiner.
\item $\beta = 1{,}6$ com \emph{gate}; \emph{hide-low} com \emph{gate}.
\item Ground truth de exposição (rostos nunca detectados, trocas de ID); decisão final do detector.
\item Equidade demográfica; avaliação humana com gravação consentida.
\end{itemize}

\section{Conclusão}
\begin{itemize}
\item Substituir rostos com identidade sintética fixa por track, empurrada para longe da identidade real, reduz a reidentificação facial automática e mantém expressão e pose melhor que censura e GANs.
\item É desidentificação facial parcial (3--9$\times$ o acaso por vídeo), não anonimização; o ganho sobre o BLANKET liberado é cobertura com menos reidentificação por rosto mostrado, não privacidade maior que censura no vídeo final.
\item Não é mais privado que blur; é outro ponto do compromisso.
\item Emoção não é preservada por nenhum método.
\item A privacidade de um método depende de quais rostos são medidos; medir o vídeo final.
\item Custo: \result{e5/pooled/adapt-0.08/class-hours} GPU-horas por hora de aula só no laço por rosto (O4 refutado); viável para estudos de poucas horas de vídeo, não para acervos de um semestre.
\end{itemize}
